In this appendix we explain the conventions used for \SU{2} and show the details of the isomorphism between \SO{4} and a class of equivalence of $\SU{2} \times \SU{2}$.


\subsection{Conventions}

We parameterise \SU{2} matrices $U$ with a vector $\vb{n} \in \R^3$ such that:
\begin{equation}
  U(\vb{n})
  =
  \cos(2 \pi n)\, \1_2
  +
  i\, \frac{\vb{n} \cdot \vb{\sigma}}{n}\, \sin(2 \pi n),
  \label{eq:su2parametrisation}
\end{equation}
where $n = \norm{\vb{n}}$ and $0 \le n \le \frac{1}{2}$.
We also identify all $\vb{n}$ when $n=\frac{1}{2}$ since in this case $U(\vb{n})= -\1_2$.
The parametrisation is such that:
\begin{eqnarray}
  U^*(\vb{n})
  & = &
  \sigma^2\, U(\vb{n})\, \sigma^2
  =
  U(\widetilde{\vb{n}}),
  \\
  U^{\dagger}(\vb{n})
  & = &
  U^T(\widetilde{\vb{n}})
  =
  U(-\vb{n}),
  \\
  -U(\vb{n})
  & = &
  U(\widehat{\vb{n}})
  \label{eq:U_props}
\end{eqnarray}
where $\sigma^2$ is the second Pauli matrix, $\widetilde{\vb{n}} = \qty( -n^1, n^2, -n^3 )$ and $\widehat{\vb{n}} = - \qty(\frac{1}{2} -n )\, \frac{\vb{n}}{n}$.

The group product of two elements $U(\vb{n} \circ \vb{m} ) = U(\vb{n})\,  U(\vb{m})$ has an explicit realisation as:
\begin{equation}
  \begin{split}
    \cos(2 \pi \norm{\vb{n} \circ \vb{m}})
    & =
    \cos(2 \pi n)\, \cos(2 \pi m)
    -
    \sin(2 \pi n)\, \sin(2\pi m)\, \frac{\vb{n} \cdot \vb{m}}{n\, m},
    \\
    \sin(2 \pi \norm{\vb{n} \circ \vb{m}})\,
    \frac{\vb{n} \circ \vb{m}}{\norm{\vb{n} \circ \vb{m}}}
    & =
    \cos(2 \pi n)\, \sin(2\pi m)\, \frac{\vb{m}}{m}
    +
    \sin(2 \pi n)\, \cos(2\pi m)\, \frac{\vb{n}}{n}.
  \end{split}
  \label{eq:product_in_SU2}
\end{equation}

\subsection{The Isomorphism}

Let $I = 1,\, 2,\, 3,\, 4$ and define:
\begin{equation}
  \tau_I = \qty( i\, \1_2,\, \vb{\sigma} ),
\end{equation}
where $\vb{\sigma} = \qty( \sigma^1,\, \sigma^2,\, \sigma^3 )$ are the Pauli matrices.
It is possible to show that:
\begin{equation}
  \begin{split}
    \qty( \tau_I )^{\dagger}
    & =
    \eta_{IJ}\, {\tau}^I,
    \\
    \qty( \tau^I )^*
    & =
    -\sigma_2\, \tau_I\, \sigma_2,
  \end{split}
  \label{eq:tau_props}
\end{equation}
where $\eta_{IJ} = \mathrm{diag}(-1,1,1,1)$.
The following relations are then a natural consequence:
\begin{eqnarray}
  \tr(\tau_I)
  & = &
  2\, i\, \delta_{I1},
  \\
  \tr(\tau_I \tau_J)
  & = &
  2\, \eta_{IJ},
  \\
  \tr(\tau_I \qty( \tau_J )^{\dagger})
  & = &
  2\, \delta_{IJ}.
\end{eqnarray}

Now consider a vector in the spinor representation:
\begin{equation}
  X_{(s)} = X^I\, \tau_I.
\end{equation}
We can recover the components using the previous properties:
\begin{equation}
  X^I
  =
  \frac{1}{2}\, \delta^{IJ}\,
  \tr(X_{(s)} \qty( \tau_J )^{\dagger})
  =
  \frac{1}{2}\, \eta^{IJ}\, \tr(X_{(s)} \tau_J),
\end{equation}
where the trace acts on the space of the $\tau$ matrices.
If the vector $X^I$ is real, using~\eqref{eq:tau_props} we have:
\begin{equation}
  \begin{split}
    X_{(s)}^{\dagger}
    & =
    X^I\, \eta_{IJ}\, \tau^J
    =
    \frac{1}{2} \tr(X_{(s)} \tau_I)\, \tau^I,
    \\
    X_{(s)}^*
    & =
    - \sigma_2\, X_{(s)}\, \sigma_2.
  \end{split}
  \label{eq:X_dagger}
\end{equation}

A rotation in spinor representation is defined as:
\begin{equation}
  X'_{(s)} = U_{L}(\vb{n})\, X_{(s)}\, U_{R}^{\dagger}(\vb{m})
\end{equation}
and it is equivalent to:
\begin{equation}
  \qty( X' )^I
  =
  \tensor{R}{^I_J}\,
  X^J
\end{equation}
through
\begin{equation}
  R_{IJ}
  =
  \frac{1}{2}
  \tr(
    \qty( \tau_I )^{\dagger}\,
    U_{L}(\vb{n})\,
    \tau_J\,
    U_{R}^{\dagger}(\vb{m})
  ).
\end{equation}
The matrix $R$ is the $4$-dimensional rotation matrix we are looking for since:
\begin{equation}
  \tr(X'_{(s)}\, (X')^{\dagger}_{(s)})
  =
  \tr(X_{(s)}\, X^{\dagger}_{(s)})
  \qquad
  \Rightarrow
  \qquad
  \finitesum{K}{1}{4} R_{IK} R^*_{JK} = \delta_{I \,J}.
\end{equation}
From the second equation in \eqref{eq:tau_props} and the first equation in \eqref{eq:U_props} we then get the reality condition on $R$:
\begin{equation}
  R_{NM}
  =
  \frac{1}{2}\, \eta_{NI}\, \eta_{MJ}\,
  \tr(\tau_I ^{\dagger}\, U_{R}\, \tau_J\, U_{L}^{\dagger})
  =
  \frac{1}{2}
  \tr(\tau_N\, U_{R}\, \tau_M^\dagger\, U_{L}^{\dagger})
  =
  R_{NM}^*.
\end{equation}
Furthermore the direct computation of the determinant of $R$ using the parametrisation~\eqref{eq:su2parametrisation} shows that $\det R = 1$.
Finally the explicit choice of the basis $\tau$ ensures $R$ to be a real matrix which ensures $R \in \SO{4}$.
Since $\qty{ U_{L},\, U_{R} }$ and $\qty{ -U_{L},\, -U_{R} }$ generate the same \SO{4} matrix then the correct isomorphism takes the form:
\begin{equation}
  \SO{4}
  \cong
  \frac{\SU{2} \times \SU{2}}{\Z_2}.
\end{equation}

